\documentclass[11pt,a4paper]{article}
\usepackage[utf8]{inputenc}
\usepackage[T1]{fontenc}
\usepackage[spanish,es-minimal,es-noshorthands]{babel}
\AtBeginDocument{\renewcommand{\sin}{\operatorname{sen}}\renewcommand{\arcsin}{\operatorname{arcsen}}}
\usepackage{lmodern}
\usepackage[sfdefault,lining]{FiraSans}
\usepackage[a4paper,top=2.3cm,bottom=2.2cm,left=1.8cm,right=1.8cm,headheight=1.3cm,headsep=0.4cm,footskip=1cm]{geometry}
\usepackage{amsmath,amssymb}
\usepackage{icomma}
\usepackage{graphicx}
\usepackage[table]{xcolor}
\usepackage{tikz}
\usepackage{fontawesome5}
\usepackage{ccicons}
\usepackage{fancyhdr}
\usepackage{lastpage}
\usepackage{tasks}
\usepackage{needspace}
\usepackage[most]{tcolorbox}
\usepackage{array}
\usepackage[hidelinks]{hyperref}
\AtBeginDocument{\hypersetup{pdftitle={Ficha de ejercicios de la unidad 2: Decimales y fracciones. 2º ESO},pdfsubject={Matemáticas 2º ESO},pdfkeywords={números decimales, decimales exactos y periódicos, operaciones con decimales, fracciones equivalentes, fracción generatriz, operaciones con fracciones, potencias de fracciones, operaciones combinadas, problemas con fracciones}}}

\definecolor{navy}{HTML}{1B2748}
\definecolor{accent}{HTML}{E8603A}
\definecolor{teal}{HTML}{1F9E90}
\definecolor{soft}{HTML}{EEF1F8}
\definecolor{muted}{HTML}{5D6685}

\setlength{\parindent}{0pt}
\setlength{\parskip}{3pt}
\settasks{label=\alph*), label-format=\bfseries\color{navy}, label-width=1.3em, item-indent=1.8em, column-sep=0.6em, after-item-skip=5pt, before-skip=3pt, after-skip=4pt}

% ---------- capçalera i peu ----------
\newcommand{\llicencia}{\href{https://creativecommons.org/licenses/by-nc-sa/4.0/deed.es}{\ccbyncsa}\ \ Licencia Creative Commons Atribución-NoComercial-CompartirIgual 4.0 Internacional (CC BY-NC-SA 4.0)}
\pagestyle{fancy}
\fancyhf{}
\fancyhead[L]{\raisebox{-0.35\height}{\includegraphics[height=0.95cm]{logo-gva.png}}}
\fancyhead[C]{\small\color{muted}Ficha de ejercicios · Unidad 2: Decimales y fracciones · 2\textsuperscript{o} ESO}
\fancyhead[R]{\raisebox{-0.35\height}{\includegraphics[height=0.95cm]{logo-vermellar.png}}}
\fancyfoot[L]{\scriptsize\color{muted}\llicencia}
\fancyfoot[R]{\small\color{muted}\thepage}
\renewcommand{\headrulewidth}{0.4pt}
\renewcommand{\headrule}{\color{muted!40}\hrule height \headrulewidth}
\fancypagestyle{primera}{\fancyhf{}\renewcommand{\headrulewidth}{0pt}
  \fancyfoot[L]{\scriptsize\color{muted}\llicencia}
  \fancyfoot[R]{\small\color{muted}\thepage}}

% ---------- nivells i exercicis ----------
\newcommand{\estrelles}[1]{{\color{accent}\ifcase#1\or\faStar\faStar[regular]\faStar[regular]\or\faStar\faStar\faStar[regular]\or\faStar\faStar\faStar\fi}}
\newcounter{exer}
\newcommand{\ex}[1]{\par\needspace{4\baselineskip}\medskip\stepcounter{exer}\noindent{\color{navy}\bfseries\theexer.}\ \ {\small\estrelles{#1}}\ \ }

\newcommand{\seccio}[2]{\par\needspace{10\baselineskip}\bigskip
  \begin{tcolorbox}[colback=navy,colframe=navy,coltext=white,arc=2mm,boxrule=0pt,left=3mm,right=3mm,top=1.2mm,bottom=1.2mm]
  {\large\bfseries #1\quad #2}
  \end{tcolorbox}\vspace{-2pt}}

\newcolumntype{C}[1]{>{\centering\arraybackslash}p{#1}}
% quadre de recordatori (només a la fitxa d'exercicis)
\newcommand{\recorda}[1]{\par\smallskip\begin{tcolorbox}[colback=accent!7,colframe=accent!55,arc=2mm,boxrule=0.5pt,left=2.5mm,right=2.5mm,top=1mm,bottom=1mm]\small{\color{accent}\faLightbulb[regular]}\ \textbf{Recuerda:}\ #1\end{tcolorbox}}
\newcommand{\res}[1]{#1}\newcommand{\rest}[1]{#1}

\def\fmneed{12}
% =====================================================================
% Capçalera FentMates de les fitxes: un color per a cada curs.
% Els colors són els de la pàgina inicial (LEVELS a assets/js/materials.js);
% el logo ve de _eines/logo/pdf/ (python3 _eines/logo/genera.py).
%
% Ús, al final del preàmbul de la unitat:
%   % =====================================================================
% Capçalera FentMates de les fitxes: un color per a cada curs.
% Els colors són els de la pàgina inicial (LEVELS a assets/js/materials.js);
% el logo ve de _eines/logo/pdf/ (python3 _eines/logo/genera.py).
%
% Ús, al final del preàmbul de la unitat:
%   % =====================================================================
% Capçalera FentMates de les fitxes: un color per a cada curs.
% Els colors són els de la pàgina inicial (LEVELS a assets/js/materials.js);
% el logo ve de _eines/logo/pdf/ (python3 _eines/logo/genera.py).
%
% Ús, al final del preàmbul de la unitat:
%   \input{../comu/fentmates.tex}
%   \fmcurs{1eso}
% i, al començament del document:
%   \capcalera{Matemàtiques}{1r ESO}{Fitxa d'exercicis · Unitat 2}{Divisibilitat}
% Opcional abans de \input: \def\fmneed{14} (espai mínim abans de cada secció).
% =====================================================================
\makeatletter
\def\fm@logodir{../../logo/pdf}
\providecommand{\fmneed}{10}
\def\fm@cursnom{}

% clau del curs, color i color del text damunt del color
\newcommand{\fm@curs}[3]{\@namedef{fm@color@#1}{#2}\@namedef{fm@text@#1}{#3}}
\fm@curs{1eso}{E6B84A}{navy}
\fm@curs{2eso}{3B82F6}{white}
\fm@curs{3eso}{FF7A50}{white}
\fm@curs{4eso-a}{A78BFA}{white}
\fm@curs{4eso-b}{6D28D9}{white}
\fm@curs{1batx-mi}{9B1C1C}{white}
\fm@curs{1batx-ccss}{EF6B6B}{white}
\fm@curs{2batx-mii}{1E6B47}{white}
\fm@curs{2batx-ccss}{2BB3A3}{white}

\newcommand{\fmcurs}[1]{%
  \@ifundefined{fm@color@#1}{\PackageError{fentmates}{Curs desconegut: #1}{}}{}%
  \edef\fm@tmp{\noexpand\definecolor{curs}{HTML}{\@nameuse{fm@color@#1}}%
               \noexpand\colorlet{curstext}{\@nameuse{fm@text@#1}}}\fm@tmp
  \colorlet{cursfosc}{curs!72!black}%
  \def\fm@logo{\fm@logodir/vertical-#1.pdf}%
  \def\fm@marca{\fm@logodir/marca-#1.pdf}}

% ---------- capçalera de la primera pàgina ----------
% #1 matèria (petit), #2 curs, #3 tipus de document i unitat, #4 títol
\newcommand{\capcalera}[4]{%
  \gdef\fm@cursnom{#2}%
  \par\noindent
  \begin{tcolorbox}[colback=curs,colframe=curs,boxrule=0pt,arc=2.2mm,width=3.5cm,height=2.15cm,
      valign=center,left=3mm,right=2mm,top=1mm,bottom=1mm,nobeforeafter,box align=center]
    \color{curstext}{\footnotesize #1}\par\vspace{1pt}{\fontsize{19}{21}\selectfont\bfseries #2\par}%
  \end{tcolorbox}\hspace{5mm}%
  \begin{minipage}[c]{\dimexpr\linewidth-3.5cm-5mm-2.9cm\relax}
    {\color{cursfosc}\small\bfseries\MakeUppercase{#3}}\par\vspace{2pt}
    {\color{navy}\fontsize{20}{23}\selectfont\bfseries #4\par}
  \end{minipage}\hfill
  \begin{minipage}[c]{2.6cm}\raggedleft\includegraphics[width=2.6cm]{\fm@logo}\end{minipage}\par
  \vspace{2pt}}

% ---------- barres de secció del color del curs ----------
\renewcommand{\seccio}[2]{\par\needspace{\fmneed\baselineskip}\bigskip
  \begin{tcolorbox}[colback=curs,colframe=curs,coltext=curstext,arc=2mm,boxrule=0pt,left=3mm,right=3mm,top=1.2mm,bottom=1.2mm]
  {\large\bfseries #1\quad #2}
  \end{tcolorbox}\vspace{-2pt}}

% ---------- estrelles de dificultat del color del curs (a Batxillerat no se'n posen) ----------
\@ifundefined{estrelles}{}{%
  \renewcommand{\estrelles}[1]{{\color{curs}\ifcase#1\or\faStar\faStar[regular]\faStar[regular]\or\faStar\faStar\faStar[regular]\or\faStar\faStar\faStar\fi}}}

% ---------- capçalera i peu de les pàgines ----------
% primera pàgina: logos institucionals al peu
\fancypagestyle{primera}{\fancyhf{}\renewcommand{\headrulewidth}{0pt}
  \fancyfoot[L]{\raisebox{-0.32\height}{\includegraphics[height=0.78cm]{logo-gva.png}}\hspace{5mm}%
    \raisebox{-0.32\height}{\includegraphics[height=0.78cm]{logo-vermellar.png}}}
  \fancyfoot[R]{\scriptsize\color{muted}\href{https://creativecommons.org/licenses/by-nc-sa/4.0/deed.ca}{\ccbyncsa}\ \ CC BY-NC-SA 4.0 · fentmates.org\qquad\small\thepage}}
% la resta: el símbol a l'esquerra i el curs a la dreta, amb el títol al centre com abans
\fancyhead[L]{\raisebox{-0.3\height}{\includegraphics[height=0.62cm]{\fm@marca}}}
\fancyhead[R]{\small\bfseries\color{cursfosc}\fm@cursnom}
\renewcommand{\headrule}{\color{curs!45}\hrule height \headrulewidth}
\makeatother

%   \fmcurs{1eso}
% i, al començament del document:
%   \capcalera{Matemàtiques}{1r ESO}{Fitxa d'exercicis · Unitat 2}{Divisibilitat}
% Opcional abans de \input: \def\fmneed{14} (espai mínim abans de cada secció).
% =====================================================================
\makeatletter
\def\fm@logodir{../../logo/pdf}
\providecommand{\fmneed}{10}
\def\fm@cursnom{}

% clau del curs, color i color del text damunt del color
\newcommand{\fm@curs}[3]{\@namedef{fm@color@#1}{#2}\@namedef{fm@text@#1}{#3}}
\fm@curs{1eso}{E6B84A}{navy}
\fm@curs{2eso}{3B82F6}{white}
\fm@curs{3eso}{FF7A50}{white}
\fm@curs{4eso-a}{A78BFA}{white}
\fm@curs{4eso-b}{6D28D9}{white}
\fm@curs{1batx-mi}{9B1C1C}{white}
\fm@curs{1batx-ccss}{EF6B6B}{white}
\fm@curs{2batx-mii}{1E6B47}{white}
\fm@curs{2batx-ccss}{2BB3A3}{white}

\newcommand{\fmcurs}[1]{%
  \@ifundefined{fm@color@#1}{\PackageError{fentmates}{Curs desconegut: #1}{}}{}%
  \edef\fm@tmp{\noexpand\definecolor{curs}{HTML}{\@nameuse{fm@color@#1}}%
               \noexpand\colorlet{curstext}{\@nameuse{fm@text@#1}}}\fm@tmp
  \colorlet{cursfosc}{curs!72!black}%
  \def\fm@logo{\fm@logodir/vertical-#1.pdf}%
  \def\fm@marca{\fm@logodir/marca-#1.pdf}}

% ---------- capçalera de la primera pàgina ----------
% #1 matèria (petit), #2 curs, #3 tipus de document i unitat, #4 títol
\newcommand{\capcalera}[4]{%
  \gdef\fm@cursnom{#2}%
  \par\noindent
  \begin{tcolorbox}[colback=curs,colframe=curs,boxrule=0pt,arc=2.2mm,width=3.5cm,height=2.15cm,
      valign=center,left=3mm,right=2mm,top=1mm,bottom=1mm,nobeforeafter,box align=center]
    \color{curstext}{\footnotesize #1}\par\vspace{1pt}{\fontsize{19}{21}\selectfont\bfseries #2\par}%
  \end{tcolorbox}\hspace{5mm}%
  \begin{minipage}[c]{\dimexpr\linewidth-3.5cm-5mm-2.9cm\relax}
    {\color{cursfosc}\small\bfseries\MakeUppercase{#3}}\par\vspace{2pt}
    {\color{navy}\fontsize{20}{23}\selectfont\bfseries #4\par}
  \end{minipage}\hfill
  \begin{minipage}[c]{2.6cm}\raggedleft\includegraphics[width=2.6cm]{\fm@logo}\end{minipage}\par
  \vspace{2pt}}

% ---------- barres de secció del color del curs ----------
\renewcommand{\seccio}[2]{\par\needspace{\fmneed\baselineskip}\bigskip
  \begin{tcolorbox}[colback=curs,colframe=curs,coltext=curstext,arc=2mm,boxrule=0pt,left=3mm,right=3mm,top=1.2mm,bottom=1.2mm]
  {\large\bfseries #1\quad #2}
  \end{tcolorbox}\vspace{-2pt}}

% ---------- estrelles de dificultat del color del curs (a Batxillerat no se'n posen) ----------
\@ifundefined{estrelles}{}{%
  \renewcommand{\estrelles}[1]{{\color{curs}\ifcase#1\or\faStar\faStar[regular]\faStar[regular]\or\faStar\faStar\faStar[regular]\or\faStar\faStar\faStar\fi}}}

% ---------- capçalera i peu de les pàgines ----------
% primera pàgina: logos institucionals al peu
\fancypagestyle{primera}{\fancyhf{}\renewcommand{\headrulewidth}{0pt}
  \fancyfoot[L]{\raisebox{-0.32\height}{\includegraphics[height=0.78cm]{logo-gva.png}}\hspace{5mm}%
    \raisebox{-0.32\height}{\includegraphics[height=0.78cm]{logo-vermellar.png}}}
  \fancyfoot[R]{\scriptsize\color{muted}\href{https://creativecommons.org/licenses/by-nc-sa/4.0/deed.ca}{\ccbyncsa}\ \ CC BY-NC-SA 4.0 · fentmates.org\qquad\small\thepage}}
% la resta: el símbol a l'esquerra i el curs a la dreta, amb el títol al centre com abans
\fancyhead[L]{\raisebox{-0.3\height}{\includegraphics[height=0.62cm]{\fm@marca}}}
\fancyhead[R]{\small\bfseries\color{cursfosc}\fm@cursnom}
\renewcommand{\headrule}{\color{curs!45}\hrule height \headrulewidth}
\makeatother

%   \fmcurs{1eso}
% i, al començament del document:
%   \capcalera{Matemàtiques}{1r ESO}{Fitxa d'exercicis · Unitat 2}{Divisibilitat}
% Opcional abans de \input: \def\fmneed{14} (espai mínim abans de cada secció).
% =====================================================================
\makeatletter
\def\fm@logodir{../../logo/pdf}
\providecommand{\fmneed}{10}
\def\fm@cursnom{}

% clau del curs, color i color del text damunt del color
\newcommand{\fm@curs}[3]{\@namedef{fm@color@#1}{#2}\@namedef{fm@text@#1}{#3}}
\fm@curs{1eso}{E6B84A}{navy}
\fm@curs{2eso}{3B82F6}{white}
\fm@curs{3eso}{FF7A50}{white}
\fm@curs{4eso-a}{A78BFA}{white}
\fm@curs{4eso-b}{6D28D9}{white}
\fm@curs{1batx-mi}{9B1C1C}{white}
\fm@curs{1batx-ccss}{EF6B6B}{white}
\fm@curs{2batx-mii}{1E6B47}{white}
\fm@curs{2batx-ccss}{2BB3A3}{white}

\newcommand{\fmcurs}[1]{%
  \@ifundefined{fm@color@#1}{\PackageError{fentmates}{Curs desconegut: #1}{}}{}%
  \edef\fm@tmp{\noexpand\definecolor{curs}{HTML}{\@nameuse{fm@color@#1}}%
               \noexpand\colorlet{curstext}{\@nameuse{fm@text@#1}}}\fm@tmp
  \colorlet{cursfosc}{curs!72!black}%
  \def\fm@logo{\fm@logodir/vertical-#1.pdf}%
  \def\fm@marca{\fm@logodir/marca-#1.pdf}}

% ---------- capçalera de la primera pàgina ----------
% #1 matèria (petit), #2 curs, #3 tipus de document i unitat, #4 títol
\newcommand{\capcalera}[4]{%
  \gdef\fm@cursnom{#2}%
  \par\noindent
  \begin{tcolorbox}[colback=curs,colframe=curs,boxrule=0pt,arc=2.2mm,width=3.5cm,height=2.15cm,
      valign=center,left=3mm,right=2mm,top=1mm,bottom=1mm,nobeforeafter,box align=center]
    \color{curstext}{\footnotesize #1}\par\vspace{1pt}{\fontsize{19}{21}\selectfont\bfseries #2\par}%
  \end{tcolorbox}\hspace{5mm}%
  \begin{minipage}[c]{\dimexpr\linewidth-3.5cm-5mm-2.9cm\relax}
    {\color{cursfosc}\small\bfseries\MakeUppercase{#3}}\par\vspace{2pt}
    {\color{navy}\fontsize{20}{23}\selectfont\bfseries #4\par}
  \end{minipage}\hfill
  \begin{minipage}[c]{2.6cm}\raggedleft\includegraphics[width=2.6cm]{\fm@logo}\end{minipage}\par
  \vspace{2pt}}

% ---------- barres de secció del color del curs ----------
\renewcommand{\seccio}[2]{\par\needspace{\fmneed\baselineskip}\bigskip
  \begin{tcolorbox}[colback=curs,colframe=curs,coltext=curstext,arc=2mm,boxrule=0pt,left=3mm,right=3mm,top=1.2mm,bottom=1.2mm]
  {\large\bfseries #1\quad #2}
  \end{tcolorbox}\vspace{-2pt}}

% ---------- estrelles de dificultat del color del curs (a Batxillerat no se'n posen) ----------
\@ifundefined{estrelles}{}{%
  \renewcommand{\estrelles}[1]{{\color{curs}\ifcase#1\or\faStar\faStar[regular]\faStar[regular]\or\faStar\faStar\faStar[regular]\or\faStar\faStar\faStar\fi}}}

% ---------- capçalera i peu de les pàgines ----------
% primera pàgina: logos institucionals al peu
\fancypagestyle{primera}{\fancyhf{}\renewcommand{\headrulewidth}{0pt}
  \fancyfoot[L]{\raisebox{-0.32\height}{\includegraphics[height=0.78cm]{logo-gva.png}}\hspace{5mm}%
    \raisebox{-0.32\height}{\includegraphics[height=0.78cm]{logo-vermellar.png}}}
  \fancyfoot[R]{\scriptsize\color{muted}\href{https://creativecommons.org/licenses/by-nc-sa/4.0/deed.ca}{\ccbyncsa}\ \ CC BY-NC-SA 4.0 · fentmates.org\qquad\small\thepage}}
% la resta: el símbol a l'esquerra i el curs a la dreta, amb el títol al centre com abans
\fancyhead[L]{\raisebox{-0.3\height}{\includegraphics[height=0.62cm]{\fm@marca}}}
\fancyhead[R]{\small\bfseries\color{cursfosc}\fm@cursnom}
\renewcommand{\headrule}{\color{curs!45}\hrule height \headrulewidth}
\makeatother

\fmcurs{2eso}

\begin{document}
\thispagestyle{primera}

\capcalera{Matemáticas}{2\textsuperscript{o} ESO}{Ficha de ejercicios · Unidad 2}{Decimales y fracciones}
\vspace{10pt}
Nombre y apellidos: \rule{6.6cm}{0.4pt}\hfill Grupo: \rule{1.6cm}{0.4pt}\hfill Fecha: \rule{2.2cm}{0.4pt}
\vspace{8pt}

\begin{tcolorbox}[colback=soft,colframe=soft,arc=2mm,left=3mm,right=3mm,top=1.5mm,bottom=1.5mm]
\small Ejercicios graduados por niveles de dificultad:\quad \estrelles{1}\ básico \qquad \estrelles{2}\ medio \qquad \estrelles{3}\ avanzado.\\[2pt] Resuélvelos en el cuaderno mostrando todos los pasos y sin calculadora. Da las fracciones en forma irreducible. Las secciones siguen el orden de los apuntes de la unidad (2.1--2.11).
\end{tcolorbox}

% =====================================================
\seccio{1}{Números decimales}

\ex{1} Escribe cómo se leen estos números:\begin{tasks}(3)
  \task $7{,}25$
  \task $0{,}408$
  \task $12{,}06$
  \task $3{,}5$
  \task $100{,}001$
  \task $0{,}09$
\end{tasks}\recorda{leemos la parte entera, la palabra \emph{unidades} y la parte decimal con el nombre de la \textbf{última} cifra: $3{,}045$ son 3 unidades y 45 milésimas.}

\ex{1} ¿Cuánto vale la cifra indicada en cada número?\begin{tasks}(3)
  \task El 7 en $34{,}571$
  \task El 3 en $2{,}738$
  \task El 7 en $607{,}15$
  \task El 8 en $0{,}0809$
  \task El 4 en $15{,}4$
  \task El 6 en $9{,}206$
\end{tasks}

\ex{1} Escribe $<$, $>$ o $=$:\begin{tasks}(3)
  \task $4{,}7\ \square\ 4{,}68$
  \task $0{,}09\ \square\ 0{,}1$
  \task $5{,}307\ \square\ 5{,}37$
  \task $12{,}5\ \square\ 12{,}5$
  \task $3{,}099\ \square\ 3{,}1$
  \task $0{,}45\ \square\ 0{,}405$
\end{tasks}

\ex{2} Ordena de menor a mayor: $2{,}5;\ 2{,}05;\ 2{,}55;\ 2{,}505;\ 2{,}15;\ 2{,}501$

\ex{2} Completa la tabla:\par\smallskip{\centering\renewcommand{\arraystretch}{1.35}\begin{tabular}{|c|C{2.4cm}|C{2.4cm}|C{2.4cm}|}\hline Número & Truncado a las centésimas & Redondeado a las décimas & Redondeado a las centésimas\\\hline $3{,}4682$ & & & \\\hline $15{,}8962$ & & & \\\hline $0{,}7049$ & & & \\\hline $29{,}996$ & & & \\\hline\end{tabular}\par}\recorda{para \textbf{redondear}, miramos la primera cifra que eliminamos: si es 0, 1, 2, 3 o 4, la última cifra se queda igual; si es 5, 6, 7, 8 o 9, le sumamos 1.}

\ex{3} Escribe tres números comprendidos entre $4{,}7$ y $4{,}71$. ¿Cuál es el número más pequeño con dos cifras decimales que, redondeado a las unidades, da 8? ¿Y el mayor?


% =====================================================
\seccio{2}{Tipos de números decimales}

\ex{1} Clasifica los números siguientes y, si son periódicos, escríbelos con la raya del periodo:\begin{tasks}(2)
  \task $4{,}252525\ldots$
  \task $0{,}125$
  \task $7{,}31666\ldots$
  \task $1{,}010010001\ldots$
  \task $0{,}777\ldots$
  \task $12{,}408408\ldots$
  \task $3{,}1252525\ldots$
  \task $\pi=3{,}14159\ldots$
\end{tasks}\recorda{periódico \textbf{puro}: el periodo empieza justo después de la coma; \textbf{mixto}: antes del periodo hay cifras que no se repiten (el anteperiodo).}

\ex{1} Indica la parte entera, el anteperiodo y el periodo de cada número:\begin{tasks}(4)
  \task $5{,}\overline{6}$
  \task $0{,}0\overline{9}$
  \task $2{,}13\overline{4}$
  \task $0{,}\overline{142857}$
\end{tasks}

\ex{2} Ordena de menor a mayor: $0{,}\overline{3};\ \ 0{,}3;\ \ 0{,}3\overline{4};\ \ 0{,}\overline{34};\ \ 0{,}33$

\ex{3} Escribe un número decimal periódico puro, uno periódico mixto, uno exacto y uno no exacto y no periódico que estén todos entre $2{,}1$ y $2{,}2$.


% =====================================================
\seccio{3}{Sumar, restar y multiplicar decimales}

\ex{1} Calcula:\begin{tasks}(3)
  \task $23{,}45+7{,}8$
  \task $40-12{,}35$
  \task $0{,}9+0{,}35$
  \task $15-6{,}47$
  \task $3{,}06+12{,}9+0{,}4$
  \task $100-0{,}01$
\end{tasks}\recorda{para sumar y restar, colocamos las \textbf{comas alineadas} y completamos con ceros.}

\ex{1} Calcula:\begin{tasks}(3)
  \task $4{,}6\cdot 3$
  \task $2{,}35\cdot 0{,}4$
  \task $0{,}6\cdot 0{,}05$
  \task $12{,}5\cdot 1{,}2$
  \task $3{,}25\cdot 1{,}4$
  \task $0{,}07\cdot 0{,}3$
\end{tasks}\recorda{multiplicamos como si no hubiera comas y separamos tantos decimales como \textbf{suman} los dos factores.}

\ex{1} Calcula mentalmente:\begin{tasks}(4)
  \task $3{,}7\cdot 100$
  \task $0{,}052\cdot 1000$
  \task $45{,}2:1000$
  \task $8:100$
  \task $6{,}15\cdot 10$
  \task $0{,}3:10$
  \task $250\cdot 0{,}1$
  \task $1{,}4\cdot 0{,}01$
\end{tasks}

\ex{2} Compro $2{,}5$ kg de manzanas a $1{,}86$ €/kg y $0{,}75$ kg de queso a $14{,}20$ €/kg. Pago con un billete de 20 €. ¿Cuánto me devuelven?

\ex{2} Una hoja de papel tiene $0{,}012$ cm de grosor. ¿Cuánto mide un paquete de 500 hojas? ¿Y una pila de 1000 paquetes, en metros?

\ex{3} Halla el número que falta: \ a) $3{,}75+\square=10$ \quad b) $\square-2{,}08=5{,}3$ \quad c) $\square\cdot100=4{,}5$ \quad d) $0{,}4\cdot\square=0{,}12$


% =====================================================
\seccio{4}{Dividir con decimales}

\ex{1} Calcula:\begin{tasks}(3)
  \task $8{,}4:4$
  \task $15:8$
  \task $9{,}75:3$
  \task $4{,}5:0{,}9$
  \task $2{,}73:0{,}7$
  \task $6:0{,}12$
\end{tasks}\recorda{si el divisor tiene decimales, multiplicamos dividendo y divisor por 10, 100, 1000… hasta que el divisor sea un número natural: $4{,}5:0{,}9=45:9$.}

\ex{2} Calcula el cociente con las cifras decimales indicadas (redondea la última):\begin{tasks}(4)
  \task $10:3$ \ (2)
  \task $7:6$ \ (2)
  \task $2:7$ \ (3)
  \task $5{,}3:0{,}9$ \ (2)
\end{tasks}

\ex{2} Calcula respetando la jerarquía de las operaciones:\begin{tasks}[after-item-skip=6pt](2)
  \task $5{,}2-1{,}5\cdot 2{,}4+0{,}6:0{,}2$
  \task $\left(3{,}5-1{,}25\right)\cdot 0{,}4+2{,}7:0{,}09$
  \task $1{,}2\cdot \left(4-2{,}5\right)+0{,}8:4$
  \task $10-{\left(0{,}5+1{,}5\right)}^{2}\cdot 2$
\end{tasks}

\ex{2} Un trozo de queso de $0{,}85$ kg cuesta $11{,}39$ €. ¿Cuál es el precio del kilo? ¿Cuánto costaría un trozo de $1{,}2$ kg?

\ex{3} Un coche gasta $31{,}5$ L de gasolina en 450 km. ¿Cuántos litros gasta cada 100 km? Si el litro cuesta $1{,}58$ €, ¿cuánto cuesta un viaje de 260 km?

\ex{3} Con una botella de $1{,}5$ L, ¿cuántos vasos de $0{,}25$ L podemos llenar? ¿Y cuántos de $0{,}2$ L? ¿Cuánto sobra en cada caso?


% =====================================================
\seccio{5}{Fracciones}

\ex{1} Calcula:\begin{tasks}(4)
  \task $\dfrac{2}{7}$ de $63$
  \task $\dfrac{3}{4}$ de $80$
  \task $\dfrac{5}{6}$ de $108$
  \task $\dfrac{3}{8}$ de $240$
\end{tasks}

\ex{1} Simplifica hasta obtener la fracción irreducible:\begin{tasks}(6)
  \task $\dfrac{84}{126}$
  \task $\dfrac{105}{180}$
  \task $\dfrac{36}{48}$
  \task $\dfrac{75}{100}$
  \task $\dfrac{240}{360}$
  \task $\dfrac{98}{140}$
\end{tasks}\recorda{dividimos numerador y denominador por su máximo común divisor (MCD).}

\ex{2} Comprueba si son equivalentes:\begin{tasks}(4)
  \task $\dfrac{14}{21}\ \text{y}\ \dfrac{22}{33}$
  \task $\dfrac{3}{4}\ \text{y}\ \dfrac{12}{15}$
  \task $\dfrac{-6}{8}\ \text{y}\ \dfrac{9}{-12}$
  \task $\dfrac{5}{7}\ \text{y}\ \dfrac{15}{20}$
\end{tasks}

\ex{2} Halla $x$ para que las fracciones sean equivalentes:\begin{tasks}(4)
  \task $\dfrac{x}{12}=\dfrac{15}{36}$
  \task $\dfrac{8}{x}=\dfrac{-2}{5}$
  \task $\dfrac{6}{9}=\dfrac{x}{15}$
  \task $\dfrac{3}{5}=\dfrac{27}{x}$
\end{tasks}

\ex{2} Ordena de menor a mayor: $\dfrac{5}{6};\ \ -\dfrac{3}{4};\ \ \dfrac{7}{9};\ \ -\dfrac{2}{3};\ \ \dfrac{3}{4}$

\ex{3} Escribe tres fracciones comprendidas entre $\dfrac{1}{3}$ y $\dfrac{1}{2}$. ¿Qué fracción de denominador 20 está justo en la mitad entre $\dfrac{2}{5}$ y $\dfrac{1}{2}$?


% =====================================================
\seccio{6}{Fracciones y decimales}

\ex{1} Sin hacer la división, di si cada fracción da un decimal exacto, periódico puro o periódico mixto. Después, compruébalo dividiendo:\begin{tasks}(4)
  \task $\dfrac{9}{20}$
  \task $\dfrac{5}{9}$
  \task $\dfrac{7}{30}$
  \task $\dfrac{21}{28}$
  \task $\dfrac{4}{11}$
  \task $\dfrac{7}{12}$
  \task $\dfrac{13}{40}$
  \task $\dfrac{1}{6}$
\end{tasks}\recorda{mira el denominador de la fracción \textbf{irreducible}: si solo tiene los factores 2 y 5, es exacto; si no tiene ninguno, periódico puro; si tiene alguno y también otros, periódico mixto.}

\ex{1} Escribe como fracción irreducible:\begin{tasks}(6)
  \task $0{,}8$
  \task $0{,}45$
  \task $1{,}2$
  \task $0{,}004$
  \task $2{,}35$
  \task $0{,}125$
\end{tasks}

\ex{2} Calcula la fracción generatriz y simplifícala:\begin{tasks}(4)
  \task $0{,}\overline{6}$
  \task $3{,}\overline{12}$
  \task $0{,}8\overline{3}$
  \task $2{,}1\overline{45}$
  \task $5{,}\overline{4}$
  \task $0{,}\overline{27}$
  \task $0{,}1\overline{6}$
  \task $1{,}2\overline{3}$
\end{tasks}\recorda{numerador: el número sin coma hasta el final del periodo menos la parte que no se repite; denominador: tantos 9 como cifras tiene el periodo y tantos 0 como cifras tiene el anteperiodo. \ $1{,}2\overline{3}=\dfrac{123-12}{90}$}

\ex{2} Comprueba, dividiendo, que $\dfrac{37}{30}=1{,}2\overline{3}$. Después, demuéstralo multiplicando $x=1{,}2\overline{3}$ por 100 y por 10 y restando.

\ex{2} Ordena de menor a mayor pasándolo todo a decimal: \ $0{,}7;\ \ \dfrac{2}{3};\ \ 0{,}67;\ \ \dfrac{3}{4};\ \ 0{,}6\overline{8}$

\ex{3} Calcula pasando a fracción: \ a) $0{,}\overline{3}+0{,}\overline{6}$ \quad b) $0{,}\overline{9}$ \quad c) $1{,}\overline{3}\cdot0{,}75$ \quad d) $0{,}1\overline{6}:0{,}\overline{5}$. ¿Qué observas en a) y b)?


% =====================================================
\seccio{7}{Suma y resta de fracciones}

\ex{1} Calcula y simplifica:\begin{tasks}[after-item-skip=6pt](3)
  \task $\dfrac{4}{7}+\dfrac{8}{7}$
  \task $\dfrac{3}{10}+\dfrac{2}{5}$
  \task $\dfrac{5}{6}-\dfrac{1}{4}$
  \task $\dfrac{7}{12}+\dfrac{5}{18}$
  \task $3-\dfrac{5}{4}$
  \task $1+\dfrac{2}{7}$
\end{tasks}\recorda{si los denominadores son distintos, reducimos a común denominador con el mcm: $\dfrac{5}{6}-\dfrac{1}{4}=\dfrac{10}{12}-\dfrac{3}{12}$.}

\ex{2} Calcula y simplifica:\begin{tasks}[after-item-skip=6pt](3)
  \task $\dfrac{1}{2}-\dfrac{2}{3}+\dfrac{1}{6}$
  \task $-\dfrac{3}{4}-\dfrac{1}{6}+2$
  \task $\dfrac{5}{6}-\dfrac{3}{4}+\dfrac{1}{3}$
  \task $\dfrac{2}{5}-\dfrac{7}{10}-1$
  \task $-\dfrac{1}{2}-\dfrac{1}{3}-\dfrac{1}{4}$
  \task $\dfrac{3}{8}-\dfrac{5}{12}+\dfrac{1}{6}$
\end{tasks}

\ex{2} Calcula (primero los paréntesis):\begin{tasks}[after-item-skip=8pt](2)
  \task $\dfrac{2}{5}-\left(\dfrac{1}{2}-\dfrac{3}{4}\right)-1$
  \task $\dfrac{2}{3}-\left[\dfrac{1}{4}-\left(\dfrac{5}{6}-1\right)\right]$
  \task $\left(1-\dfrac{1}{3}\right)-\left(\dfrac{1}{2}-\dfrac{1}{5}\right)$
  \task $\dfrac{3}{4}-\left[2-\left(\dfrac{1}{3}+\dfrac{5}{6}\right)\right]$
\end{tasks}

\ex{3} ¿Qué fracción hay que sumar a $\dfrac{3}{8}$ para obtener $\dfrac{5}{6}$? ¿Cuál hay que restar a $\dfrac{1}{4}$ para obtener $-\dfrac{2}{3}$?

\ex{3} Encuentra el error y corrígelo: \ $\dfrac{2}{3}+\dfrac{1}{5}=\dfrac{2+1}{3+5}=\dfrac{3}{8}$


% =====================================================
\seccio{8}{Multiplicación y división de fracciones}

\ex{1} Calcula y simplifica:\begin{tasks}[after-item-skip=6pt](3)
  \task $\dfrac{2}{3}\cdot \dfrac{5}{7}$
  \task $\dfrac{3}{2}\cdot \dfrac{2}{5}$
  \task $\dfrac{3}{7}\cdot 14$
  \task $\dfrac{4}{3}:\dfrac{5}{6}$
  \task $\dfrac{3}{5}:\dfrac{9}{10}$
  \task $6:\dfrac{3}{4}$
\end{tasks}\recorda{para dividir, multiplicamos por la \textbf{inversa}: $\dfrac{a}{b}:\dfrac{c}{d}=\dfrac{a}{b}\cdot\dfrac{d}{c}=\dfrac{a\cdot d}{b\cdot c}$.}

\ex{1} Escribe la inversa de cada número:\begin{tasks}(5)
  \task $\dfrac{2}{5}$
  \task $-\dfrac{7}{4}$
  \task $3$
  \task $-\dfrac{1}{6}$
  \task $1$
\end{tasks}

\ex{2} Calcula y simplifica:\begin{tasks}[after-item-skip=6pt](3)
  \task $\left(-\dfrac{8}{15}\right)\cdot \left(-\dfrac{25}{12}\right)$
  \task $\left(-\dfrac{5}{6}\right):\dfrac{10}{9}$
  \task $\dfrac{7}{10}:\left(-14\right)$
  \task $\dfrac{14}{15}\cdot \dfrac{25}{21}$
  \task $\dfrac{3}{4}\cdot \dfrac{8}{9}:\dfrac{2}{3}$
  \task $\dfrac{2}{5}:\left(\dfrac{4}{5}\cdot \dfrac{3}{8}\right)$
\end{tasks}

\ex{2} Calcula: \ a) $\dfrac{2}{3}$ de $\dfrac{9}{10}$ \quad b) la mitad de $\dfrac{4}{5}$ \quad c) el triple de $\dfrac{5}{12}$ \quad d) $\dfrac{\;\frac{3}{4}\;}{\frac{9}{2}}$

\ex{3} ¿Por qué fracción hay que multiplicar $\dfrac{4}{9}$ para obtener $\dfrac{2}{3}$? ¿Entre qué fracción hay que dividir $\dfrac{5}{8}$ para obtener 5?


% =====================================================
\seccio{9}{Potencias de fracciones}

\ex{1} Calcula:\begin{tasks}[after-item-skip=6pt](4)
  \task ${\left(\dfrac{3}{4}\right)}^{3}$
  \task ${\left(-\dfrac{2}{5}\right)}^{2}$
  \task ${\left(-\dfrac{1}{3}\right)}^{3}$
  \task ${\left(\dfrac{7}{9}\right)}^{0}$
  \task ${\left(\dfrac{1}{2}\right)}^{5}$
  \task ${\left(-\dfrac{3}{2}\right)}^{4}$
  \task ${\left(-\dfrac{2}{3}\right)}^{1}$
  \task ${\left(\dfrac{5}{2}\right)}^{2}$
\end{tasks}\recorda{$\left(\dfrac{a}{b}\right)^n=\dfrac{a^n}{b^n}$. Con base negativa, el resultado es positivo si el exponente es par y negativo si es impar.}

\ex{1} Calcula y compara: \ $\left(-\dfrac{2}{3}\right)^2$, \ $-\left(\dfrac{2}{3}\right)^2$, \ $\dfrac{2^2}{3}$ \ y \ $\dfrac{-2^2}{3^2}$.

\ex{2} Escribe como una sola potencia y calcula:\begin{tasks}[after-item-skip=8pt](2)
  \task ${\left(\dfrac{2}{3}\right)}^{5}\cdot {\left(\dfrac{2}{3}\right)}^{2}:{\left(\dfrac{2}{3}\right)}^{4}$
  \task ${\left(\dfrac{3}{5}\right)}^{7}:{\left(\dfrac{3}{5}\right)}^{5}$
  \task ${\left[{\left(-\dfrac{1}{2}\right)}^{3}\right]}^{2}:{\left(-\dfrac{1}{2}\right)}^{4}$
  \task ${\left(\dfrac{4}{7}\right)}^{6}:{\left(\dfrac{4}{7}\right)}^{5}$
  \task ${\left[{\left(\dfrac{1}{3}\right)}^{2}\right]}^{2}$
  \task ${\left(\dfrac{5}{4}\right)}^{3}\cdot {\left(\dfrac{5}{4}\right)}^{0}$
\end{tasks}\recorda{con la misma base: $x^m\cdot x^n=x^{m+n}$, \ $x^m:x^n=x^{m-n}$, \ $(x^m)^n=x^{m\cdot n}$.}

\ex{2} Calcula de la manera más rápida, usando las propiedades: \ a) $\left(\dfrac{5}{6}\right)^4\cdot\left(\dfrac{3}{10}\right)^4\cdot4^4$ \quad b) $\left(\dfrac{9}{4}\right)^3:\left(\dfrac{3}{2}\right)^3$

\ex{3} Simplifica: \ $\dfrac{\left(\frac{2}{5}\right)^3\cdot\left(\frac{2}{5}\right)^4}{\left[\left(\frac{2}{5}\right)^2\right]^3}$ \qquad y \qquad $\dfrac{\left(-\frac{3}{4}\right)^5\cdot\left(-\frac{3}{4}\right)}{\left(-\frac{3}{4}\right)^3\cdot\left(-\frac{3}{4}\right)^2}$


% =====================================================
\seccio{10}{Operaciones combinadas}

\ex{1} Calcula respetando la jerarquía de las operaciones:\begin{tasks}[after-item-skip=8pt](3)
  \task $\dfrac{1}{2}+\dfrac{1}{3}\cdot \dfrac{3}{4}$
  \task $\left(\dfrac{1}{2}+\dfrac{1}{3}\right)\cdot \dfrac{3}{4}$
  \task $1-\dfrac{2}{3}:\dfrac{4}{9}$
  \task $\dfrac{3}{5}+\dfrac{2}{5}\cdot \dfrac{10}{3}$
  \task $\dfrac{5}{2}+\dfrac{1}{3}\cdot \dfrac{4}{5}$
  \task $\dfrac{4}{3}\cdot \left(\dfrac{2}{5}+\dfrac{1}{3}\right)$
\end{tasks}\recorda{orden: 1) paréntesis; 2) potencias; 3) multiplicaciones y divisiones, de izquierda a derecha; 4) sumas y restas.}

\ex{2} Calcula:\begin{tasks}[after-item-skip=8pt](2)
  \task $\dfrac{1}{2}-{\left(\dfrac{2}{3}\right)}^{2}\cdot \dfrac{3}{4}$
  \task ${\left(\dfrac{2}{3}\right)}^{2}-\dfrac{1}{3}\cdot \left(2-\dfrac{1}{2}\right)$
  \task $\dfrac{5}{6}:\left(\dfrac{1}{2}-\dfrac{1}{3}\right)-{2}^{2}$
  \task ${\left(-\dfrac{1}{2}\right)}^{3}+\dfrac{3}{4}:\left(1-\dfrac{1}{3}\right)$
  \task ${\left(\dfrac{3}{4}-1\right)}^{2}\cdot 8+\dfrac{1}{2}$
  \task $\dfrac{21}{4}-\dfrac{2}{5}:\dfrac{4}{3}$
\end{tasks}

\ex{2} Calcula:\begin{tasks}[after-item-skip=8pt](2)
  \task $2-\left[\dfrac{1}{3}+{\left(\dfrac{1}{2}\right)}^{2}\right]:\dfrac{5}{6}$
  \task ${\left[\dfrac{3}{4}-{\left(\dfrac{1}{2}\right)}^{2}\right]}^{2}:{\left(-\dfrac{1}{2}\right)}^{3}+\dfrac{5}{3}\cdot \dfrac{6}{5}$
  \task $\left(\dfrac{1}{4}+\dfrac{2}{3}-\dfrac{1}{12}\right)\cdot 3+\left(\dfrac{3}{2}-\dfrac{5}{4}\right):\dfrac{1}{2}$
  \task $\dfrac{3}{2}:\dfrac{1}{4}\cdot \dfrac{1}{5}+\dfrac{3}{2}\cdot \dfrac{7}{5}$
\end{tasks}

\ex{3} Calcula: \ $\dfrac{\frac{2}{3}-\frac{1}{4}}{1+\frac{1}{6}}-\left(\dfrac{3}{4}-1\right)^2\cdot8$

\ex{3} Pon los paréntesis necesarios para que cada igualdad sea cierta: \ a) $\dfrac{1}{2}+\dfrac{1}{2}\cdot\dfrac{2}{3}=\dfrac{2}{3}$ \quad b) $1-\dfrac{1}{3}:\dfrac{1}{3}+1=\dfrac{1}{2}$


% =====================================================
\seccio{11}{Problemas con fracciones}

\ex{1} De los 28 alumnos de una clase, $\dfrac{3}{7}$ son chicos. ¿Cuántos chicos y cuántas chicas hay?

\ex{1} He leído 120 páginas de un libro, que son los $\dfrac{3}{5}$ del total. ¿Cuántas páginas tiene el libro? ¿Cuántas me quedan?

\ex{1} Un bocadillo cuesta $3{,}60$ € y el pan representa $\dfrac{1}{4}$; el jamón, $\dfrac{2}{3}$. ¿Cuánto cuesta el pan? ¿Y el jamón? ¿Cuánto cuesta el resto?

\ex{2} De toda una cena, Joan se ha comido $\dfrac{1}{7}$, Elena $\dfrac{2}{5}$ y Maite $\dfrac{1}{8}$. ¿Qué fracción se han comido entre los tres? ¿Qué fracción ha sobrado? ¿Quién ha comido más?

\ex{2} Un depósito de 240 L está lleno. El lunes se gasta $\dfrac{1}{4}$ y el martes, $\dfrac{2}{3}$ de lo que quedaba. ¿Cuántos litros quedan? ¿Qué fracción del depósito es?

\ex{2} En una votación, $\dfrac{3}{8}$ de los votantes votan sí, $\dfrac{1}{3}$ votan no y el resto, 14 personas, se abstienen. ¿Cuántas personas han votado?

\ex{2} Una botella llena pesa $1{,}2$ kg y, cuando hemos vaciado $\dfrac{3}{4}$, pesa $0{,}48$ kg. ¿Cuánto pesa la botella vacía?

\ex{3} Marta gasta $\dfrac{2}{5}$ de la paga en ropa y $\dfrac{1}{3}$ de lo que le queda en el cine. Todavía le quedan 24 €. ¿Cuánto cobraba de paga?

\ex{3} Un ciclista hace $\dfrac{2}{7}$ de una etapa por la mañana y $\dfrac{3}{5}$ de lo que le queda por la tarde. Todavía le faltan 40 km. ¿Cuántos kilómetros tiene la etapa? ¿Cuántos ha hecho por la mañana?

\ex{3} Un padre reparte 180 € entre sus tres hijos: al mayor le da $\dfrac{2}{5}$, al mediano $\dfrac{5}{8}$ de lo que queda y al pequeño, el resto. ¿Cuánto recibe cada uno? ¿Qué fracción del total recibe el pequeño?

\ex{3} En una pastelería hacen 150 cruasanes. Por la mañana venden $\dfrac{3}{5}$ y por la tarde, $\dfrac{3}{4}$ de los que quedan. Los que sobran se venden a mitad de precio, a $0{,}60$ € cada uno. ¿Cuántos sobran y cuánto se recauda por ellos?

\vfill
\begin{center}\color{muted}\small\faCheckCircle[regular]\ \ ¡Ánimo! Comprueba tus resultados con las herramientas interactivas de los apuntes de la unidad.\end{center}

\end{document}
